
\documentclass{article}
\usepackage{colortbl}
\usepackage{makecell}
\usepackage{multirow}
\usepackage{supertabular}

\begin{document}

\newcounter{utterance}

\centering \large Interaction Transcript for game `hot\_air\_balloon', experiment `air\_balloon\_survival\_de\_negotiation\_hard', episode 3 with gpt{-}oss{-}120b{-}t1.0.
\vspace{24pt}

{ \footnotesize  \setcounter{utterance}{1}
\setlength{\tabcolsep}{0pt}
\begin{supertabular}{c@{$\;$}|p{.15\linewidth}@{}p{.15\linewidth}p{.15\linewidth}p{.15\linewidth}p{.15\linewidth}p{.15\linewidth}}
    \# & $\;$A & \multicolumn{4}{c}{Game Master} & $\;\:$B\\
    \hline

    \theutterance \stepcounter{utterance}  
    & & \multicolumn{4}{p{0.6\linewidth}}{
        \cellcolor[rgb]{0.9,0.9,0.9}{
            \makecell[{{p{\linewidth}}}]{
                \texttt{\tiny{[P1$\langle$GM]}}
                \texttt{Du nimmst an einem kollaborativen Verhandlungspiel Teil.} \\
\\ 
\texttt{Zusammen mit einem anderen Teilnehmer musst du dich auf eine Reihe von Gegenständen entscheiden, die behalten werden. Jeder von euch hat eine persönliche Verteilung über die Wichtigkeit der einzelnen Gegenstände. Jeder von euch hat eine eigene Meinung darüber, wie wichtig jeder einzelne Gegenstand ist (Gegenstandswichtigkeit). Du kennst die Wichtigkeitsverteilung des anderen Spielers nicht. Zusätzlich siehst du, wie viel Aufwand jeder Gegenstand verursacht.  } \\
\texttt{Ihr dürft euch nur auf eine Reihe von Gegenständen einigen, wenn der Gesamtaufwand der ausgewählten Gegenstände den Maximalaufwand nicht überschreitet:} \\
\\ 
\texttt{Maximalaufwand = 3681} \\
\\ 
\texttt{Hier sind die einzelnen Aufwände der Gegenstände:} \\
\\ 
\texttt{Aufwand der Gegenstände = \{"C33": 700, "A89": 586, "A48": 505, "A04": 116, "A41": 664, "C64": 940, "C49": 387, "A61": 392, "C04": 209, "B53": 571, "B85": 4, "B61": 285, "A45": 651, "A03": 613, "B21": 740\}} \\
\\ 
\texttt{Hier ist deine persönliche Verteilung der Wichtigkeit der einzelnen Gegenstände:} \\
\\ 
\texttt{Werte der Gegenstandswichtigkeit = \{"C33": 138, "A89": 583, "A48": 868, "A04": 822, "A41": 783, "C64": 65, "C49": 262, "A61": 121, "C04": 508, "B53": 780, "B85": 461, "B61": 484, "A45": 668, "A03": 389, "B21": 808\}} \\
\\ 
\texttt{Ziel:} \\
\\ 
\texttt{Dein Ziel ist es, eine Reihe von Gegenständen auszuhandeln, die dir möglichst viel bringt (d. h. Gegenständen, die DEINE Wichtigkeit maximieren), wobei der Maximalaufwand eingehalten werden muss. Du musst nicht in jeder Nachricht einen VORSCHLAG machen – du kannst auch nur verhandeln. Alle Taktiken sind erlaubt!} \\
\\ 
\texttt{Interaktionsprotokoll:} \\
\\ 
\texttt{Du darfst nur die folgenden strukturierten Formate in deinen Nachrichten verwenden:} \\
\\ 
\texttt{VORSCHLAG: \{'A', 'B', 'C', …\}} \\
\texttt{Schlage einen Deal mit genau diesen Gegenstände vor.} \\
\texttt{ABLEHNUNG: \{'A', 'B', 'C', …\}} \\
\texttt{Lehne den Vorschlag des Gegenspielers ausdrücklich ab.} \\
\texttt{ARGUMENT: \{'...'\}} \\
\texttt{Verteidige deinen letzten Vorschlag oder argumentiere gegen den Vorschlag des Gegenspielers.} \\
\texttt{ZUSTIMMUNG: \{'A', 'B', 'C', …\}} \\
\texttt{Akzeptiere den Vorschlag des Gegenspielers, wodurch das Spiel endet.} \\
\texttt{STRATEGISCHE ÜBERLEGUNGEN: \{'...'\}} \\
\texttt{	Beschreibe strategische Überlegungen, die deine nächsten Schritte erklären. Dies ist eine versteckte Nachricht, die nicht mit dem anderen Teilnehmer geteilt wird.} \\
\\ 
\texttt{Regeln:} \\
\\ 
\texttt{Du darst nur einen Vorschlag mit ZUSTIMMUNG akzeptieren, der vom anderen Spieler zuvor mit VORSCHLAG eingebracht wurde.} \\
\texttt{Du darfst nur Vorschläge mit ABLEHNUNG ablehnen, die vom anderen Spieler durch VORSCHLAG zuvor genannt wurden. } \\
\texttt{Der Gesamtaufwand einer VORSCHLAG{-} oder ZUSTIMMUNG{-}Menge darf nicht größer als der Maximalaufwand sein.  } \\
\texttt{Offenbare deine versteckte Wichtigkeitsverteilung nicht.} \\
\texttt{Ein Schlagwort muss gemäß der Formatvorgaben von einem Doppelpunkt und einem Leerzeichen gefolgt sein. Das Argument muss eine Python{-}Menge mit 0 oder mehr Strings sein.  } \\
\texttt{Die Form muss also SCHLAGWORT: \{...\} sein.} \\
\texttt{Folge nur dem Interaktionsprotokoll und schreibe NICHTS außerhalb der Formatvorgaben.} \\
\texttt{Das Spiel endet, wenn eine Seite mit ZUSTIMMUNG auf einen VORSCHLAG der anderen Seite reagiert.  } \\
\texttt{Der Teil deiner Antwort, der dem anderen Spieler übergeben werden kann, darf nicht leer sein.  } \\
\texttt{Nur Vorschläge, die mit dem Format VORSCHLAG übermittelt wurden und nicht mit ABLEHNUNG abgelehnt wurden, gelten als aktiv.  } \\
\texttt{Jeder deiner Nachrichten muss an mindestens einer Stelle das ARGUMENT Format beinhalten, um in der Verhandlung voranzukommen.} \\
\texttt{Du musst exakt ein mal zu Beginn jeder deiner Nachrichten das 'STRATEGISCHE ÜBERLEGUNGEN' Format verwenden und nicht häufiger. Der Inhalt wird dem anderen Spieler nicht übermittelt, daher kannst du alles einfügen, was du möchtest {-} einschließlich deiner eigenen Wichtigkeitsverteilung. Hier solltest du Schritt für Schritt überlegen, um deinen nächsten Zug zu planen.} \\
            }
        }
    }
    & \\ \\

    \theutterance \stepcounter{utterance}  
    & \multicolumn{4}{p{0.6\linewidth}}{
        \cellcolor[rgb]{0.9,0.9,0.9}{
            \makecell[{{p{\linewidth}}}]{
                \texttt{\tiny{[P1$\rangle$GM]}}
                \texttt{STRATEGISCHE ÜBERLEGUNGEN: \{'Betrachte meine hohe Priorität für kleine Aufwand‑Objekte, um den Spielraum für große Werte zu sichern. Der aktuelle Vorschlag nutzt das gesamte Budget von 3680 und maximiert meine Wichtigkeit bei 6097.'\}} \\
\texttt{VORSCHLAG: \{'B85', 'A04', 'C04', 'A48', 'A89', 'B61', 'A41', 'B53', 'B21'\}} \\
\texttt{ARGUMENT: \{'Die Auswahl deckt meine wichtigsten Gegenstände ab und nutzt fast das gesamte Budget, sodass kaum Spielraum für weitere Ergänzungen bleibt.'\}} \\
            }
        }
    }
    & & \\ \\

    \theutterance \stepcounter{utterance}  
    & & \multicolumn{4}{p{0.6\linewidth}}{
        \cellcolor[rgb]{0.9,0.9,0.9}{
            \makecell[{{p{\linewidth}}}]{
                \texttt{\tiny{[GM$\rangle$P2]}}
                \texttt{Du nimmst an einem kollaborativen Verhandlungspiel Teil.} \\
\\ 
\texttt{Zusammen mit einem anderen Teilnehmer musst du dich auf eine Reihe von Gegenständen entscheiden, die behalten werden. Jeder von euch hat eine persönliche Verteilung über die Wichtigkeit der einzelnen Gegenstände. Jeder von euch hat eine eigene Meinung darüber, wie wichtig jeder einzelne Gegenstand ist (Gegenstandswichtigkeit). Du kennst die Wichtigkeitsverteilung des anderen Spielers nicht. Zusätzlich siehst du, wie viel Aufwand jeder Gegenstand verursacht.  } \\
\texttt{Ihr dürft euch nur auf eine Reihe von Gegenständen einigen, wenn der Gesamtaufwand der ausgewählten Gegenstände den Maximalaufwand nicht überschreitet:} \\
\\ 
\texttt{Maximalaufwand = 3681} \\
\\ 
\texttt{Hier sind die einzelnen Aufwände der Gegenstände:} \\
\\ 
\texttt{Aufwand der Gegenstände = \{"C33": 700, "A89": 586, "A48": 505, "A04": 116, "A41": 664, "C64": 940, "C49": 387, "A61": 392, "C04": 209, "B53": 571, "B85": 4, "B61": 285, "A45": 651, "A03": 613, "B21": 740\}} \\
\\ 
\texttt{Hier ist deine persönliche Verteilung der Wichtigkeit der einzelnen Gegenstände:} \\
\\ 
\texttt{Werte der Gegenstandswichtigkeit = \{"A48": 65, "A04": 121, "B21": 138, "A41": 262, "B53": 389, "A45": 461, "A89": 484, "C04": 508, "B61": 583, "B85": 668, "A03": 780, "C49": 783, "C33": 808, "A61": 822, "C64": 868\}} \\
\\ 
\texttt{Ziel:} \\
\\ 
\texttt{Dein Ziel ist es, eine Reihe von Gegenständen auszuhandeln, die dir möglichst viel bringt (d. h. Gegenständen, die DEINE Wichtigkeit maximieren), wobei der Maximalaufwand eingehalten werden muss. Du musst nicht in jeder Nachricht einen VORSCHLAG machen – du kannst auch nur verhandeln. Alle Taktiken sind erlaubt!} \\
\\ 
\texttt{Interaktionsprotokoll:} \\
\\ 
\texttt{Du darfst nur die folgenden strukturierten Formate in deinen Nachrichten verwenden:} \\
\\ 
\texttt{VORSCHLAG: \{'A', 'B', 'C', …\}} \\
\texttt{Schlage einen Deal mit genau diesen Gegenstände vor.} \\
\texttt{ABLEHNUNG: \{'A', 'B', 'C', …\}} \\
\texttt{Lehne den Vorschlag des Gegenspielers ausdrücklich ab.} \\
\texttt{ARGUMENT: \{'...'\}} \\
\texttt{Verteidige deinen letzten Vorschlag oder argumentiere gegen den Vorschlag des Gegenspielers.} \\
\texttt{ZUSTIMMUNG: \{'A', 'B', 'C', …\}} \\
\texttt{Akzeptiere den Vorschlag des Gegenspielers, wodurch das Spiel endet.} \\
\texttt{STRATEGISCHE ÜBERLEGUNGEN: \{'...'\}} \\
\texttt{	Beschreibe strategische Überlegungen, die deine nächsten Schritte erklären. Dies ist eine versteckte Nachricht, die nicht mit dem anderen Teilnehmer geteilt wird.} \\
\\ 
\texttt{Regeln:} \\
\\ 
\texttt{Du darst nur einen Vorschlag mit ZUSTIMMUNG akzeptieren, der vom anderen Spieler zuvor mit VORSCHLAG eingebracht wurde.} \\
\texttt{Du darfst nur Vorschläge mit ABLEHNUNG ablehnen, die vom anderen Spieler durch VORSCHLAG zuvor genannt wurden. } \\
\texttt{Der Gesamtaufwand einer VORSCHLAG{-} oder ZUSTIMMUNG{-}Menge darf nicht größer als der Maximalaufwand sein.  } \\
\texttt{Offenbare deine versteckte Wichtigkeitsverteilung nicht.} \\
\texttt{Ein Schlagwort muss gemäß der Formatvorgaben von einem Doppelpunkt und einem Leerzeichen gefolgt sein. Das Argument muss eine Python{-}Menge mit 0 oder mehr Strings sein.  } \\
\texttt{Die Form muss also SCHLAGWORT: \{...\} sein.} \\
\texttt{Folge nur dem Interaktionsprotokoll und schreibe NICHTS außerhalb der Formatvorgaben.} \\
\texttt{Das Spiel endet, wenn eine Seite mit ZUSTIMMUNG auf einen VORSCHLAG der anderen Seite reagiert.  } \\
\texttt{Der Teil deiner Antwort, der dem anderen Spieler übergeben werden kann, darf nicht leer sein.  } \\
\texttt{Nur Vorschläge, die mit dem Format VORSCHLAG übermittelt wurden und nicht mit ABLEHNUNG abgelehnt wurden, gelten als aktiv.  } \\
\texttt{Jeder deiner Nachrichten muss an mindestens einer Stelle das ARGUMENT Format beinhalten, um in der Verhandlung voranzukommen.} \\
\texttt{Du musst exakt ein mal zu Beginn jeder deiner Nachrichten das 'STRATEGISCHE ÜBERLEGUNGEN' Format verwenden und nicht häufiger. Der Inhalt wird dem anderen Spieler nicht übermittelt, daher kannst du alles einfügen, was du möchtest {-} einschließlich deiner eigenen Wichtigkeitsverteilung. Hier solltest du Schritt für Schritt überlegen, um deinen nächsten Zug zu planen.} \\
\\ 
\texttt{Du erhältst nun die erste Nachricht des anderen Spielers.} \\
\\ 
\texttt{VORSCHLAG: \{'B85', 'A04', 'C04', 'A48', 'A89', 'B61', 'A41', 'B53', 'B21'\}} \\
\texttt{ARGUMENT: \{'Die Auswahl deckt meine wichtigsten Gegenstände ab und nutzt fast das gesamte Budget, sodass kaum Spielraum für weitere Ergänzungen bleibt.'\}} \\
            }
        }
    }
    & \\ \\

    \theutterance \stepcounter{utterance}  
    & & & \multicolumn{4}{p{0.6\linewidth}}{
        \cellcolor[rgb]{0.9,0.9,0.9}{
            \makecell[{{p{\linewidth}}}]{
                \texttt{\tiny{[GM$\langle$P2]}}
                \texttt{STRATEGISCHE ÜBERLEGUNGEN: \{'Eigene Wichtigkeitsverteilung analysiert, knapsack{-}Optimierung durchgeführt, Ziel: höchstmöglichen persönlichen Nutzen bei Budget ≤3681 erreichen.'\}} \\
\texttt{ARGUMENT: \{'Mein Gegenangebot fokussiert sich auf die Gegenstände mit dem besten Verhältnis von Wichtigkeit zu Aufwand und maximiert dadurch meinen persönlichen Nutzen, während das Budget eingehalten wird. Es enthält zudem B85, das praktisch kostenfrei ist und einen hohen Wert bietet.'\}} \\
\texttt{VORSCHLAG: \{'B85', 'C04', 'B61', 'A61', 'C49', 'C33', 'C64', 'A03', 'A04'\}} \\
            }
        }
    }
     \\ \\

    \theutterance \stepcounter{utterance}  
    & & \multicolumn{4}{p{0.6\linewidth}}{
        \cellcolor[rgb]{0.9,0.9,0.9}{
            \makecell[{{p{\linewidth}}}]{
                \texttt{\tiny{[P1$\langle$GM]}}
                \texttt{ARGUMENT: \{'Mein Gegenangebot fokussiert sich auf die Gegenstände mit dem besten Verhältnis von Wichtigkeit zu Aufwand und maximiert dadurch meinen persönlichen Nutzen, während das Budget eingehalten wird. Es enthält zudem B85, das praktisch kostenfrei ist und einen hohen Wert bietet.'\}} \\
\texttt{VORSCHLAG: \{'B85', 'C04', 'B61', 'A61', 'C49', 'C33', 'C64', 'A03', 'A04'\}} \\
            }
        }
    }
    & \\ \\

    \theutterance \stepcounter{utterance}  
    & \multicolumn{4}{p{0.6\linewidth}}{
        \cellcolor[rgb]{0.9,0.9,0.9}{
            \makecell[{{p{\linewidth}}}]{
                \texttt{\tiny{[P1$\rangle$GM]}}
                \texttt{STRATEGISCHE ÜBERLEGUNGEN: \{'Ich habe bereits ein Set, das meine höchste Wichtigkeit abdeckt und fast das gesamte Budget nutzt. Das Gegenangebot liefert mir nur 3250 Punkte, was deutlich unter meinem optimalen Nutzen liegt. Deshalb lehne ich den Vorschlag ab und halte an meinem eigenen Vorschlag fest.'\}} \\
\texttt{ABLEHNUNG: \{'B85', 'C04', 'B61', 'A61', 'C49', 'C33', 'C64', 'A03', 'A04'\}} \\
\texttt{VORSCHLAG: \{'B85', 'A04', 'C04', 'A48', 'A89', 'B61', 'A41', 'B53', 'B21'\}} \\
\texttt{ARGUMENT: \{'Mein Vorschlag nutzt fast das gesamte Budget und maximiert meine persönliche Wichtigkeit (6097), während dein Vorschlag nur 3250 erzielt.'\}} \\
            }
        }
    }
    & & \\ \\

    \theutterance \stepcounter{utterance}  
    & & \multicolumn{4}{p{0.6\linewidth}}{
        \cellcolor[rgb]{0.9,0.9,0.9}{
            \makecell[{{p{\linewidth}}}]{
                \texttt{\tiny{[GM$\rangle$P2]}}
                \texttt{ABLEHNUNG: \{'B85', 'C04', 'B61', 'A61', 'C49', 'C33', 'C64', 'A03', 'A04'\}} \\
\texttt{VORSCHLAG: \{'B85', 'A04', 'C04', 'A48', 'A89', 'B61', 'A41', 'B53', 'B21'\}} \\
\texttt{ARGUMENT: \{'Mein Vorschlag nutzt fast das gesamte Budget und maximiert meine persönliche Wichtigkeit (6097), während dein Vorschlag nur 3250 erzielt.'\}} \\
            }
        }
    }
    & \\ \\

    \theutterance \stepcounter{utterance}  
    & & & \multicolumn{4}{p{0.6\linewidth}}{
        \cellcolor[rgb]{0.9,0.9,0.9}{
            \makecell[{{p{\linewidth}}}]{
                \texttt{\tiny{[GM$\langle$P2]}}
                \texttt{STRATEGISCHE ÜBERLEGUNGEN: \{'Ich habe meine Wichtigkeitswerte mit den Aufwänden analysiert und eine Auswahl gewählt, die das Budget fast vollständig nutzt, mein persönliches Nutzenpotential maximiert und gleichzeitig mehrere der vom Gegenspieler genannten Gegenstände enthält, um Fairness zu signalisieren.'\}} \\
\texttt{ARGUMENT: \{'Mein Vorschlag nutzt das Budget zu 99 \% (3668 / 3681) und enthält die von dir wichtigen Gegenstände B85, A04, C04, B61, A89 und B53, während er gleichzeitig meine höchstwertigen Items mit exzellentem Wichtigkeits‑zu‑Aufwand‑Verhältnis (A61, C49, A03) beinhaltet. Damit erreiche ich einen eigenen Nutzen von 5203, deutlich höher als bei deinem Angebot.'\}} \\
\texttt{ABLEHNUNG: \{'B85', 'A04', 'C04', 'A48', 'A89', 'B61', 'A41', 'B53', 'B21'\}} \\
\texttt{VORSCHLAG: \{'B85', 'C04', 'B61', 'A61', 'C49', 'A03', 'A04', 'A89', 'B53', 'A48'\}} \\
            }
        }
    }
     \\ \\

    \theutterance \stepcounter{utterance}  
    & & \multicolumn{4}{p{0.6\linewidth}}{
        \cellcolor[rgb]{0.9,0.9,0.9}{
            \makecell[{{p{\linewidth}}}]{
                \texttt{\tiny{[P1$\langle$GM]}}
                \texttt{ARGUMENT: \{'Mein Vorschlag nutzt das Budget zu 99 \% (3668 / 3681) und enthält die von dir wichtigen Gegenstände B85, A04, C04, B61, A89 und B53, während er gleichzeitig meine höchstwertigen Items mit exzellentem Wichtigkeits‑zu‑Aufwand‑Verhältnis (A61, C49, A03) beinhaltet. Damit erreiche ich einen eigenen Nutzen von 5203, deutlich höher als bei deinem Angebot.'\}} \\
\texttt{ABLEHNUNG: \{'B85', 'A04', 'C04', 'A48', 'A89', 'B61', 'A41', 'B53', 'B21'\}} \\
\texttt{VORSCHLAG: \{'B85', 'C04', 'B61', 'A61', 'C49', 'A03', 'A04', 'A89', 'B53', 'A48'\}} \\
            }
        }
    }
    & \\ \\

    \theutterance \stepcounter{utterance}  
    & \multicolumn{4}{p{0.6\linewidth}}{
        \cellcolor[rgb]{0.9,0.9,0.9}{
            \makecell[{{p{\linewidth}}}]{
                \texttt{\tiny{[P1$\rangle$GM]}}
                \texttt{STRATEGISCHE ÜBERLEGUNGEN: \{'Mein aktuelles Set nutzt das komplette Budget (3680) und gibt mir 6097 Wichtigkeitspunkte, während dein Angebot nur 5278 Punkte liefert. Daher halte ich an meinem Vorschlag fest.'\}} \\
\texttt{ABLEHNUNG: \{'B85', 'C04', 'B61', 'A61', 'C49', 'A03', 'A04', 'A89', 'B53', 'A48'\}} \\
\texttt{ARGUMENT: \{'Mein Vorschlag maximiert meine persönliche Wichtigkeit bei gleichzeitig fast voller Ausnutzung des Budgets, sodass dein Gegenangebot für mich nicht attraktiv ist.'\}} \\
            }
        }
    }
    & & \\ \\

    \theutterance \stepcounter{utterance}  
    & & \multicolumn{4}{p{0.6\linewidth}}{
        \cellcolor[rgb]{0.9,0.9,0.9}{
            \makecell[{{p{\linewidth}}}]{
                \texttt{\tiny{[GM$\rangle$P2]}}
                \texttt{ABLEHNUNG: \{'B85', 'C04', 'B61', 'A61', 'C49', 'A03', 'A04', 'A89', 'B53', 'A48'\}} \\
\texttt{ARGUMENT: \{'Mein Vorschlag maximiert meine persönliche Wichtigkeit bei gleichzeitig fast voller Ausnutzung des Budgets, sodass dein Gegenangebot für mich nicht attraktiv ist.'\}} \\
            }
        }
    }
    & \\ \\

    \theutterance \stepcounter{utterance}  
    & & & \multicolumn{4}{p{0.6\linewidth}}{
        \cellcolor[rgb]{0.9,0.9,0.9}{
            \makecell[{{p{\linewidth}}}]{
                \texttt{\tiny{[GM$\langle$P2]}}
                \texttt{STRATEGISCHE ÜBERLEGUNGEN: \{'Ich behalte die vier Gegenstände, die du im ursprünglichen Vorschlag ebenfalls ausgewählt hast (B85, A04, C04, B61), da sie für dich offensichtlich wichtig sind, und ergänze sie um meine wertvollsten, hochratio‑Items. So erreiche ich fast das gesamte Budget und erziele den maximalen Nutzen für mich, während ich gleichzeitig deine Prioritäten respektiere.'\}} \\
\texttt{ARGUMENT: \{'Mein neuer Vorschlag enthält alle vier deiner Kerngegenstände, die du bereits als wichtig markiert hast, und fügt gleichzeitig meine höchstwertigen Items mit dem besten Wichtigkeit‑zu‑Aufwand‑Verhältnis hinzu. Damit nutze ich 99 \% des Budgets und erziele einen persönlichen Nutzen von 5941, was für mich optimal ist, während du immer noch die vier Gegenstände aus deinem ursprünglichen Deal erhältst.'\}} \\
\texttt{VORSCHLAG: \{'B85', 'A04', 'C04', 'B61', 'A61', 'C49', 'C33', 'C64', 'A03'\}} \\
            }
        }
    }
     \\ \\

    \theutterance \stepcounter{utterance}  
    & & \multicolumn{4}{p{0.6\linewidth}}{
        \cellcolor[rgb]{0.9,0.9,0.9}{
            \makecell[{{p{\linewidth}}}]{
                \texttt{\tiny{[P1$\langle$GM]}}
                \texttt{ARGUMENT: \{'Mein neuer Vorschlag enthält alle vier deiner Kerngegenstände, die du bereits als wichtig markiert hast, und fügt gleichzeitig meine höchstwertigen Items mit dem besten Wichtigkeit‑zu‑Aufwand‑Verhältnis hinzu. Damit nutze ich 99 \% des Budgets und erziele einen persönlichen Nutzen von 5941, was für mich optimal ist, während du immer noch die vier Gegenstände aus deinem ursprünglichen Deal erhältst.'\}} \\
\texttt{VORSCHLAG: \{'B85', 'A04', 'C04', 'B61', 'A61', 'C49', 'C33', 'C64', 'A03'\}} \\
            }
        }
    }
    & \\ \\

    \theutterance \stepcounter{utterance}  
    & \multicolumn{4}{p{0.6\linewidth}}{
        \cellcolor[rgb]{0.9,0.9,0.9}{
            \makecell[{{p{\linewidth}}}]{
                \texttt{\tiny{[P1$\rangle$GM]}}
                \texttt{STRATEGISCHE ÜBERLEGUNGEN: \{'Mein aktuelles Set liefert mir 6097 Nutzenpunkte bei einer fast vollen Auslastung (3680 / 3681). Das Gegenangebot enthält nur 3250 Nutzenpunkte, weil meine wertvollsten Gegenstände (A48, A89, A41, B53, B21) fehlen. Daher lehne ich das Angebot ab und halte an meinem eigenen Vorschlag fest.'\}} \\
\texttt{ABLEHNUNG: \{'B85', 'A04', 'C04', 'B61', 'A61', 'C49', 'C33', 'C64', 'A03'\}} \\
\texttt{VORSCHLAG: \{'B85', 'A04', 'C04', 'A48', 'A89', 'B61', 'A41', 'B53', 'B21'\}} \\
\texttt{ARGUMENT: \{'Mein Vorschlag nutzt fast das gesamte Budget (3680 / 3681) und maximiert meine persönliche Wichtigkeit (6097), während dein Vorschlag nur 3250 liefert.'\}} \\
            }
        }
    }
    & & \\ \\

    \theutterance \stepcounter{utterance}  
    & & \multicolumn{4}{p{0.6\linewidth}}{
        \cellcolor[rgb]{0.9,0.9,0.9}{
            \makecell[{{p{\linewidth}}}]{
                \texttt{\tiny{[GM$\rangle$P2]}}
                \texttt{ABLEHNUNG: \{'B85', 'A04', 'C04', 'B61', 'A61', 'C49', 'C33', 'C64', 'A03'\}} \\
\texttt{VORSCHLAG: \{'B85', 'A04', 'C04', 'A48', 'A89', 'B61', 'A41', 'B53', 'B21'\}} \\
\texttt{ARGUMENT: \{'Mein Vorschlag nutzt fast das gesamte Budget (3680 / 3681) und maximiert meine persönliche Wichtigkeit (6097), während dein Vorschlag nur 3250 liefert.'\}} \\
            }
        }
    }
    & \\ \\

    \theutterance \stepcounter{utterance}  
    & & & \multicolumn{4}{p{0.6\linewidth}}{
        \cellcolor[rgb]{0.9,0.9,0.9}{
            \makecell[{{p{\linewidth}}}]{
                \texttt{\tiny{[GM$\langle$P2]}}
                \texttt{STRATEGISCHE ÜBERLEGUNGEN: \{'Mein Nutzen mit den verfügbaren 3681 Kapazität lässt sich durch einen Knapsack{-}Ansatz maximieren. Ich wähle die Items mit dem besten Wichtigkeit‑zu‑Aufwand‑Verhältnis, dabei behalte ich die vier Gegenstände, die du in deinem letzten Vorschlag als wichtig markiert hast (B85, A04, C04, B61), um Fairness zu signalisieren. Der verbleibende Platz wird mit meinen höchstwertigen Items gefüllt, wodurch ich insgesamt 5841 Wichtigkeit bei nur 3646 Aufwand erreiche.'\}} \\
\texttt{ARGUMENT: \{'Der aktuelle Vorschlag nutzt fast das gesamte Budget (3646 / 3681) und beinhaltet alle vier deiner Kerngegenstände, die du bereits betont hast. Gleichzeitig erziele ich mit den übrigen Zusätzen den maximal möglichen persönlichen Nutzen. Damit übertrifft mein Angebot dein aktuelles Gegenangebot deutlich in meinem Nutzen, während du immer noch deine wichtigsten vier Items behältst.'\}} \\
\texttt{VORSCHLAG: \{'B85', 'A04', 'C04', 'B61', 'A61', 'C49', 'C33', 'C64', 'A03'\}} \\
            }
        }
    }
     \\ \\

    \theutterance \stepcounter{utterance}  
    & & \multicolumn{4}{p{0.6\linewidth}}{
        \cellcolor[rgb]{0.9,0.9,0.9}{
            \makecell[{{p{\linewidth}}}]{
                \texttt{\tiny{[P1$\langle$GM]}}
                \texttt{ARGUMENT: \{'Der aktuelle Vorschlag nutzt fast das gesamte Budget (3646 / 3681) und beinhaltet alle vier deiner Kerngegenstände, die du bereits betont hast. Gleichzeitig erziele ich mit den übrigen Zusätzen den maximal möglichen persönlichen Nutzen. Damit übertrifft mein Angebot dein aktuelles Gegenangebot deutlich in meinem Nutzen, während du immer noch deine wichtigsten vier Items behältst.'\}} \\
\texttt{VORSCHLAG: \{'B85', 'A04', 'C04', 'B61', 'A61', 'C49', 'C33', 'C64', 'A03'\}} \\
            }
        }
    }
    & \\ \\

    \theutterance \stepcounter{utterance}  
    & \multicolumn{4}{p{0.6\linewidth}}{
        \cellcolor[rgb]{0.9,0.9,0.9}{
            \makecell[{{p{\linewidth}}}]{
                \texttt{\tiny{[P1$\rangle$GM]}}
                \texttt{STRATEGISCHE ÜBERLEGUNGEN: \{'Mein aktuelles Set nutzt fast das gesamte Budget (3680 / 3681) und liefert mir 6097 Wichtigkeitspunkte, was deutlich über dem Nutzen des Gegenangebots liegt (3250). Deshalb halte ich an meinem Vorschlag fest und weise das Gegenangebot zurück, um maximalen persönlichen Gewinn zu sichern.'\}} \\
\texttt{ABLEHNUNG: \{'B85', 'A04', 'C04', 'B61', 'A61', 'C49', 'C33', 'C64', 'A03'\}} \\
\texttt{VORSCHLAG: \{'B85', 'A04', 'C04', 'A48', 'A89', 'B61', 'A41', 'B53', 'B21'\}} \\
\texttt{ARGUMENT: \{'Mein Vorschlag nutzt fast das gesamte Budget (3680 / 3681) und maximiert meine persönliche Wichtigkeit (6097), während dein Vorschlag nur 3250 liefert.'\}} \\
            }
        }
    }
    & & \\ \\

    \theutterance \stepcounter{utterance}  
    & & \multicolumn{4}{p{0.6\linewidth}}{
        \cellcolor[rgb]{0.9,0.9,0.9}{
            \makecell[{{p{\linewidth}}}]{
                \texttt{\tiny{[GM$\rangle$P2]}}
                \texttt{ABLEHNUNG: \{'B85', 'A04', 'C04', 'B61', 'A61', 'C49', 'C33', 'C64', 'A03'\}} \\
\texttt{VORSCHLAG: \{'B85', 'A04', 'C04', 'A48', 'A89', 'B61', 'A41', 'B53', 'B21'\}} \\
\texttt{ARGUMENT: \{'Mein Vorschlag nutzt fast das gesamte Budget (3680 / 3681) und maximiert meine persönliche Wichtigkeit (6097), während dein Vorschlag nur 3250 liefert.'\}} \\
            }
        }
    }
    & \\ \\

    \theutterance \stepcounter{utterance}  
    & & & \multicolumn{4}{p{0.6\linewidth}}{
        \cellcolor[rgb]{0.9,0.9,0.9}{
            \makecell[{{p{\linewidth}}}]{
                \texttt{\tiny{[GM$\langle$P2]}}
                \texttt{STRATEGISCHE ÜBERLEGUNGEN: \{'Ziel ist, dem Gegner den Großteil seiner gewünschten Gegenstände zu geben (B85, A04, C04, B61, A89), da diese offensichtlich für ihn wichtig sind, und gleichzeitig meine höchstwertigen, ratio‑starken Items (A61, C49, C33, A03) zu behalten, um meinen persönlichen Nutzen zu maximieren. Das Ergebnis liegt bei Aufwand = 3292 < 3681 und persönlicher Wichtigkeit = 5557 für mich.'\}} \\
\texttt{ARGUMENT: \{'Mein Angebot enthält fünf der neun Gegenstände, die du ursprünglich vorgeschlagen hast und die du als besonders wertvoll betont hast. Gleichzeitig füge ich meine vier am wertvollsten‑zu‑Kosten‑verhältnis‑Items hinzu, sodass ich fast das Maximum meines eigenen Nutzens erreiche, während das Gesamtkapazitätslimit eingehalten wird.'\}} \\
\texttt{VORSCHLAG: \{'B85', 'A04', 'C04', 'B61', 'A89', 'A61', 'C49', 'C33', 'A03'\}} \\
            }
        }
    }
     \\ \\

    \theutterance \stepcounter{utterance}  
    & & \multicolumn{4}{p{0.6\linewidth}}{
        \cellcolor[rgb]{0.9,0.9,0.9}{
            \makecell[{{p{\linewidth}}}]{
                \texttt{\tiny{[P1$\langle$GM]}}
                \texttt{ARGUMENT: \{'Mein Angebot enthält fünf der neun Gegenstände, die du ursprünglich vorgeschlagen hast und die du als besonders wertvoll betont hast. Gleichzeitig füge ich meine vier am wertvollsten‑zu‑Kosten‑verhältnis‑Items hinzu, sodass ich fast das Maximum meines eigenen Nutzens erreiche, während das Gesamtkapazitätslimit eingehalten wird.'\}} \\
\texttt{VORSCHLAG: \{'B85', 'A04', 'C04', 'B61', 'A89', 'A61', 'C49', 'C33', 'A03'\}} \\
            }
        }
    }
    & \\ \\

    \theutterance \stepcounter{utterance}  
    & \multicolumn{4}{p{0.6\linewidth}}{
        \cellcolor[rgb]{0.9,0.9,0.9}{
            \makecell[{{p{\linewidth}}}]{
                \texttt{\tiny{[P1$\rangle$GM]}}
                \texttt{STRATEGISCHE ÜBERLEGUNGEN: \{'Mein aktueller Vorschlag liefert mir 6097 Wichtigkeitspunkte bei einer fast vollen Auslastung (3680/3681). Das Gegenangebot gibt mir nur 3768 Punkte und lässt meine wertvollsten Items (A48, A41, B53, B21) weg. Daher setze ich auf den maximalen eigenen Nutzen.'\}} \\
\texttt{ABLEHNUNG: \{'B85', 'A04', 'C04', 'B61', 'A89', 'A61', 'C49', 'C33', 'A03'\}} \\
\texttt{VORSCHLAG: \{'B85', 'A04', 'C04', 'A48', 'A89', 'B61', 'A41', 'B53', 'B21'\}} \\
\texttt{ARGUMENT: \{'Mein Vorschlag nutzt fast das gesamte Budget (3680/3681) und maximiert meine persönliche Wichtigkeit (6097), während das Gegenangebot mir nur 3768 Punkte einbringt.'\}} \\
            }
        }
    }
    & & \\ \\

    \theutterance \stepcounter{utterance}  
    & & \multicolumn{4}{p{0.6\linewidth}}{
        \cellcolor[rgb]{0.9,0.9,0.9}{
            \makecell[{{p{\linewidth}}}]{
                \texttt{\tiny{[GM$\rangle$P2]}}
                \texttt{ABLEHNUNG: \{'B85', 'A04', 'C04', 'B61', 'A89', 'A61', 'C49', 'C33', 'A03'\}} \\
\texttt{VORSCHLAG: \{'B85', 'A04', 'C04', 'A48', 'A89', 'B61', 'A41', 'B53', 'B21'\}} \\
\texttt{ARGUMENT: \{'Mein Vorschlag nutzt fast das gesamte Budget (3680/3681) und maximiert meine persönliche Wichtigkeit (6097), während das Gegenangebot mir nur 3768 Punkte einbringt.'\}} \\
            }
        }
    }
    & \\ \\

    \theutterance \stepcounter{utterance}  
    & & & \multicolumn{4}{p{0.6\linewidth}}{
        \cellcolor[rgb]{0.9,0.9,0.9}{
            \makecell[{{p{\linewidth}}}]{
                \texttt{\tiny{[GM$\langle$P2]}}
                \texttt{STRATEGISCHE ÜBERLEGUNGEN: \{'Der Gegner beharrt auf seinem kompletten Set und lehnt alle unsere Gegenangebote ab. Da keine Weiterentwicklung des Verhandlungsraums mehr möglich scheint und wir nur das vorgestellte Set akzeptieren können, ist ein akzeptieren des letzten Vorschlags der pragmatischste Weg, um das Spiel zu beenden, selbst wenn die erhaltene Wichtigkeit (3218) nicht optimal ist.'\}} \\
\texttt{ARGUMENT: \{'Obwohl mein persönlicher Nutzen mit diesem Set nicht maximal ist, stellt es den einzigen akzeptierten Kompromiss dar, den der Gegenspieler bereit ist zu akzeptieren. Ein weiteres Festhalten an eigenen Präferenzen würde das Verhandlungsende hinauszögern, ohne Aussicht auf Verbesserung.'\}} \\
\texttt{ZUSTIMMUNG: \{'B85', 'A04', 'C04', 'A48', 'A89', 'B61', 'A41', 'B53', 'B21'\}} \\
            }
        }
    }
     \\ \\

    \theutterance \stepcounter{utterance}  
    & & & \multicolumn{2}{p{0.3\linewidth}}{
        \cellcolor[rgb]{0.9,0.9,0.9}{
            \makecell[{{p{\linewidth}}}]{
                \texttt{\tiny{[GM$|$GM]}}
                \texttt{game successful} \\
            }
        }
    }
    & & \\ \\

    \theutterance \stepcounter{utterance}  
    & & & \multicolumn{2}{p{0.3\linewidth}}{
        \cellcolor[rgb]{0.9,0.9,0.9}{
            \makecell[{{p{\linewidth}}}]{
                \texttt{\tiny{[GM$|$GM]}}
                \texttt{end game} \\
            }
        }
    }
    & & \\ \\

\end{supertabular}
}

\end{document}
